\documentclass[a4paper,11pt]{article}
\usepackage[utf8]{inputenc}
\usepackage[T1]{fontenc}
\usepackage{geometry}
\usepackage{enumitem}
\usepackage{titlesec}
\usepackage[table]{xcolor}
\usepackage{fancyhdr}
\usepackage{tabularx}
\usepackage{listings}
\usepackage{lastpage}
\usepackage{hyperref}
\usepackage{booktabs}

% Page Setup
\geometry{top=1in, bottom=1in, left=0.8in, right=0.8in}
\pagestyle{fancy}
\fancyhf{}
\lhead{\small Project 83: Agentic AI Compiler Assistant}
\rhead{\small Week 15 Report}
\cfoot{\thepage\ of \pageref{LastPage}}

% Formatting
\titleformat{\section}{\Large\bfseries\color{cyan!70!black}}{\thesection.}{1em}{}[\titlerule]
\titleformat{\subsection}{\large\bfseries\color{cyan!50!black}}{\thesubsection}{1em}{}

% Code Listings
\lstset{
    basicstyle=\ttfamily\small, breaklines=true, frame=single,
    backgroundcolor=\color{gray!10}, keywordstyle=\color{blue},
    stringstyle=\color{green!50!black}, commentstyle=\color{gray},
    showstringspaces=false, tabsize=4, captionpos=b
}

\newcommand{\statusgood}[1]{\colorbox{green!20}{\textbf{#1}}}
\newcommand{\statuscritical}[1]{\colorbox{red!20}{\textbf{#1}}}
\newcommand{\statuswarn}[1]{\colorbox{yellow!20}{\textbf{#1}}}

\begin{document}

\begin{center}
    {\huge \textbf{Week 15 Progress Report}} \\
    \vspace{3mm}
    {\large \textbf{Project 83: Conversational Agentic AI Compiler Assistant}} \\
    \vspace{2mm}
    \textit{Documentation \& Report Finalization --- Academic Report, Slides \& Submission Package} \\
    \vspace{2mm}
    \textbf{Reporting Period:} Week 15 | \textbf{Status:} \statusgood{COMPLETED}
\end{center}

\vspace{5mm}

\section{Executive Summary}

Week 15 is the \textbf{Documentation \& Report Finalization} sprint. The entire body of project documentation --- spanning 15 weekly reports, the Software Requirements Specification (SRS), Literature Review, Architecture Diagrams, Prompt Templates, and a Final Comprehensive Report --- has been reviewed, cross-referenced, and finalized. A \textbf{15-slide presentation deck} summarizing the project for the live demonstration has been produced. The GitHub repository has been tagged as \texttt{v1.0.0} and a submission-ready archive (ZIP) has been created.

\textbf{Key Achievement:} The project documentation is now \textbf{complete and internally consistent}. All code snippets in reports match the actual source code; all test counts match the CI run results; all architecture diagrams reflect the final implemented system. The project is ready for the Week 16 live demonstration.

\subsection{Week Overview}
\begin{itemize}[leftmargin=*]
    \item \textbf{Planned Objectives:} Final report, presentation slides, README, API docs, submission package
    \item \textbf{Achieved:} All objectives; inline docstrings added across all \texttt{src/} modules
    \item \textbf{Test Coverage:} No new tests this week; all 392 tests remain passing
\end{itemize}

\section{Documentation Inventory}

\subsection{2.1 Reports Produced (Weeks 1--15)}

\begin{table}[h]
\centering
\begin{tabularx}{\textwidth}{|c|l|c|c|}
\hline
\rowcolor{cyan!10}
\textbf{Week} & \textbf{Title} & \textbf{Pages} & \textbf{Status} \\
\hline
1 & Problem Definition \& EBNF Grammar Design & 12 & \statusgood{FINAL} \\
2 & Gemini API Setup \& Agentic Workflow Design & 15 & \statusgood{FINAL} \\
3 & Requirements \& Security Policy Analysis (SRS) & 14 & \statusgood{FINAL} \\
4 & System Architecture \& Orchestration Planning & 22 & \statusgood{FINAL} \\
5 & Lexer (Lexical Analyzer) Development & 11 & \statusgood{FINAL} \\
6 & Syntax Parser (Recursive Descent) & 13 & \statusgood{FINAL} \\
7 & Semantic Analyzer \& Basic Diagnostics & 21 & \statusgood{FINAL} \\
8 & Gemini API Integration \& Context Injection & 11 & \statusgood{FINAL} \\
9 & Agentic Tool-Use (Function Calling) & 13 & \statusgood{FINAL} \\
10 & Conversational Debugging Interface & 11 & \statusgood{FINAL} \\
11 & Security Risk Auditing Agent (SAST) & 16 & \statusgood{FINAL} \\
12 & Performance Benchmarking & 14 & \statusgood{FINAL} \\
13 & Explainability \& Traceability (XAI) & 15 & \statusgood{FINAL} \\
14 & Integration Testing \& Bug Squashing & 16 & \statusgood{FINAL} \\
15 & Documentation \& Report Finalization & (this) & \statusgood{FINAL} \\
\hline
\end{tabularx}
\caption{Complete weekly report inventory}
\end{table}

\subsection{2.2 Supporting Documents}

\begin{table}[h]
\centering
\begin{tabularx}{\textwidth}{|l|l|c|}
\hline
\rowcolor{cyan!10}
\textbf{Document} & \textbf{Location} & \textbf{Status} \\
\hline
Software Requirements Specification (SRS) & \texttt{reports/SRS\_latex.tex} & \statusgood{FINAL} \\
Literature Review & \texttt{reports/literature\_review.tex} & \statusgood{FINAL} \\
Detailed Execution Plan & \texttt{docs/DetailedExecutionPlanandReport.tex} & \statusgood{FINAL} \\
Project Description \& Methodology & \texttt{docs/ProjectDescriptionMethodologicalApproach.tex} & \statusgood{FINAL} \\
Weekly Progress Consolidated Document & \texttt{docs/WeeklyProgressReport.tex} & \statusgood{FINAL} \\
Comprehensive Project Overview & \texttt{reports/comprehensive\_project\_overview.tex} & \statusgood{FINAL} \\
Final Comprehensive Report & \texttt{reports/final\_report.tex} & \statusgood{FINAL} \\
Prompt Template Library & \texttt{docs/prompts.md} & \statusgood{FINAL} \\
README.md & \texttt{README.md} & \statusgood{FINAL} \\
\hline
\end{tabularx}
\caption{Supporting document inventory}
\end{table}

\section{Final Comprehensive Report Structure}

The \texttt{final\_report.tex} (produced this week) is a single, self-contained academic report structured as follows:

\begin{enumerate}
    \item \textbf{Abstract} --- 250-word summary of objectives, methodology, and results
    \item \textbf{Introduction} --- Motivation, problem statement, scope of Mini-Python subset
    \item \textbf{Literature Review} --- Related work on compiler construction, LLM agents, SAST tools
    \item \textbf{System Architecture} --- Three-layer design (Compiler front-end, Agentic orchestration, UI)
    \item \textbf{Implementation: Compiler Front-End} (Weeks 5--7) --- Lexer, Parser, Semantic Analyzer
    \item \textbf{Implementation: AI Integration} (Weeks 8--9) --- Context Provider, ReAct agent, tool registry
    \item \textbf{Implementation: Conversational Interface} (Week 10) --- ChatSession, CLI, Streamlit UI
    \item \textbf{Implementation: Security Auditing} (Week 11) --- SAST engine, threat categories, OWASP mapping
    \item \textbf{Evaluation: Benchmarking Results} (Week 12) --- Latency, accuracy, memory tables
    \item \textbf{Evaluation: Explainability} (Week 13) --- ReasoningLog, AST visualization, XAI analysis
    \item \textbf{Integration \& Testing} (Week 14) --- E2E suite, bug inventory, CI/CD
    \item \textbf{Discussion} --- Limitations, future work, ethical considerations
    \item \textbf{Conclusion} --- Summary of contributions
    \item \textbf{References} --- 25+ academic and technical references
    \item \textbf{Appendices} --- EBNF grammar, full SAST rule table, sample prompt templates
\end{enumerate}

\section{Presentation Deck Summary}

The 15-slide presentation (\texttt{docs/Project83\_Presentation.pdf}) is structured as:

\begin{table}[h]
\centering
\begin{tabularx}{\textwidth}{|c|l|X|}
\hline
\rowcolor{cyan!10}
\textbf{Slide} & \textbf{Title} & \textbf{Key Content} \\
\hline
1 & Title & Project 83, team, date \\
2 & Problem Statement & Why compilers are not enough; beginner pain points; security gaps \\
3 & Project Goals & Three pillars: Accuracy, Intelligence, Security \\
4 & System Architecture & Three-layer diagram; data flow from source code to conversation \\
5 & Compiler Front-End & Lexer $\to$ Parser $\to$ AST $\to$ Semantic Analyzer; code example \\
6 & AI Integration & Context injection; ReAct loop; tool-use function calling \\
7 & Conversational Debugger & Multi-turn chat; CLI + Streamlit UI; session state \\
8 & Security Auditing (SAST) & Five threat categories; OWASP mapping; live demo screenshot \\
9 & Benchmarking Results & Latency table; accuracy chart (93\% vs 52\%); memory footprint \\
10 & Explainability (XAI) & ReasoningLog trace; annotated AST screenshot \\
11 & Integration \& Testing & E2E closed-loop diagram; 392 tests; CI/CD pipeline \\
12 & Live Demo Preview & Screenshot of full Streamlit UI \\
13 & Challenges \& Solutions & Top 5 key challenges with resolutions \\
14 & Future Work & Multi-language support; real LLM fine-tuning; VS Code extension \\
15 & Conclusion \& Q\&A & Summary of contributions; repository link \\
\hline
\end{tabularx}
\caption{Presentation deck structure}
\end{table}

\section{README \& API Documentation}

The \texttt{README.md} was fully rewritten this week with:
\begin{itemize}[leftmargin=*]
    \item \textbf{Badges}: CI passing, Python 3.11/3.12, 392 tests, MIT License
    \item \textbf{Architecture diagram}: ASCII art pipeline overview
    \item \textbf{Quick Start}: 5-step setup guide (clone, install, set API key, run demo, run tests)
    \item \textbf{Module Reference}: Each \texttt{src/} module listed with its public API
    \item \textbf{Demo Scripts}: All 11 demo scripts with expected output descriptions
    \item \textbf{Security Policy}: How to report vulnerabilities found via the SAST scanner
\end{itemize}

All \texttt{src/} modules now have complete NumPy-style docstrings. A \texttt{pydoc} index was generated automatically via \texttt{python -m pydoc -w src/}.

\section{Submission Package}

The GitHub repository was tagged \texttt{v1.0.0}. A ZIP submission archive was created containing:

\begin{lstlisting}[language=bash, caption=Submission archive contents]
project-83-v1.0.0/
  src/             # Full source code (compiler/ + orchestrator/ + benchmarks/)
  tests/           # 392 tests (Weeks 5-14)
  reports/         # 15 weekly reports + SRS + literature review (LaTeX)
  docs/            # Execution plan, methodology, prompts, consolidated progress
  grammar/         # EBNF grammar files
  demo_week*.py    # 12 demonstration scripts
  app_week*.py     # 3 Streamlit applications (Weeks 9, 10, 11)
  README.md
  requirements.txt
  .github/workflows/ci.yml
\end{lstlisting}

\section{Deliverables}

\begin{itemize}[leftmargin=*]
    \item \statusgood{COMPLETE} \texttt{reports/final\_report.tex} --- Final comprehensive academic report (14 sections)
    \item \statusgood{COMPLETE} \texttt{docs/WeeklyProgressReport.tex} --- Consolidated weekly progress document
    \item \statusgood{COMPLETE} \texttt{docs/Project83\_Presentation.pdf} --- 15-slide presentation deck
    \item \statusgood{COMPLETE} \texttt{README.md} --- Fully rewritten with badges, quick start, module reference
    \item \statusgood{COMPLETE} All \texttt{src/} modules --- NumPy-style docstrings added
    \item \statusgood{COMPLETE} GitHub tag \texttt{v1.0.0} and release ZIP archive
    \item \statusgood{COMPLETE} \texttt{reports/week\_15\_report.tex} --- This report
\end{itemize}

\section{Next Week Preview}

\textbf{Week 16: Live Demonstration \& Submission}

\begin{itemize}[leftmargin=*]
    \item Deliver the live demo using the Streamlit UI and CLI
    \item Record a Loom walkthrough video of the complete system
    \item Submit the final source code repository and documentation package
    \item Present to the assessment panel using the 15-slide deck
\end{itemize}

\section{Conclusion}

Week 15 finalised the entire documentation body of Project 83. With 15 weekly reports, a comprehensive final academic report, presentation slides, and a fully documented codebase, the project is ready for its Week 16 live demonstration and formal submission. All 392 tests remain green, the CI/CD pipeline is active, and the submission package is prepared.

\vspace{5mm}
\noindent\textbf{Prepared by:} Project Team \\
\textbf{Reporting Date:} Week 15 Completion (19 April 2026) \\
\textbf{Next Review:} Week 16 -- Live Demonstration \& Submission

\end{document}
